\scanpage{284}
Soit $S$ un préschéma de base, $A$ un schéma abélien sur $S$, $B$ un schéma sur $S$, $\Gamma$ un groupe fini, opérant sur $A$ par automorphismes et sur $B$ sans pt fixe, alors $\Gamma$ opère sur $A\times B$ diagonalement, et on peut considérer $X=(A\times B)/\Gamma$ comme un préschéma \deleted{sur} \uline{abélien} sur $Y=B/\Gamma$. On a donc, si \deleted{\unsure{on suppose que}} comme $X/Y$ \deleted{est} a une section,
\[
\uuline{\operatorname{Pic}}_{X/S}\simeq\uuline{\operatorname{Pic}}_{Y/S}\times\prod_{Y/S}\uuline{\operatorname{Pic}}_{X/Y}
\]
\margin{\diagram{\begin{tikzcd}[column sep=small,row sep=small]
X\arrow[d,"f"']&A\times_S B\arrow[l,"\Gamma"']\arrow[d,"\unsure{\mathrm{pr}_1}"]\\
Y\arrow[d,"g"']&B\arrow[l,"\Gamma"]\arrow[dl,no head]\\S&
\end{tikzcd}}}
donc, si \inserted{$B\longrightarrow S$ donc} $g$ est disons projectif et simple, et $A\longrightarrow S$ projectif, on aura
\[
\uline{\operatorname{Pic}}_{X/S}\simeq\uline{\operatorname{Pic}}_{Y/S}\times\prod_{Y/S}\uline{\operatorname{Pic}}_{X/Y}
\]
Or on a
\[
\prod_{Y/S}\uuline{\operatorname{Pic}}_{X/Y}\simeq
\inserted{\prod_{Y/S}(A'\times B)/\Gamma\simeq}
\uuline{\operatorname{Hom}}_{S,\Gamma}(B,\uline{\operatorname{Pic}}_{A/S})
\]
\inserted{\uline{Soit}} $A'$ est le schéma abélien dual de $A$, sur lequel on fait opérer $\Gamma$ par ``transport de structure'';
\margin{et $\uline{N}$ le schéma de Néron-\unsure{Sévéri}. On a une suite exacte
\diagram{\begin{tikzcd}[column sep=small]
0\arrow[r]&\uuline{\operatorname{Hom}}_{S,\Gamma}(B,A')\arrow[r]&\uuline{\operatorname{Hom}}_{S,\Gamma}(B,\uline{\operatorname{Pic}}_{A/S})\arrow[r]&\uuline{\operatorname{Hom}}_{S,\Gamma}(B,\uline{\operatorname{NS}}_{A/S})\arrow[d,equal]\\&&&\uuline{N}^{\Gamma}
\end{tikzcd}}
d'où
\[
\uuline{\operatorname{Pic}}_{X/S}^{\tau}\simeq\uuline{\operatorname{Pic}}_{Y/S}^{\tau}\times\uuline{\operatorname{Hom}}_{S,\Gamma}(B,A')^{\tau}
\]
}
\uline{Si} $A'^{\Gamma}$ \uline{est} \unsure{---} \uline{à fibres discrètes} sur $S$ [c'est le cas si parmi les opérations de $\Gamma$ il y a la symétrie $x\mapsto -x$] \inserted{\uline{si $B$ un} schéma abélien, sur lequel} \deleted{\unsure{le groupe} \unsure{qui}\illegible{3}} \inserted{\unsure{opère par translations}}, on voit facilement que
\[
\uuline{\operatorname{Hom}}_{S,\Gamma}(B,A')^{\inserted{\text{\textcircled{$\tau$}}}}\simeq A'^{\Gamma},
\]
d'où
\figure{\begin{tikzpicture}[x=1cm,y=1cm]
\draw (0,2.3)..controls (0,2.1) and (0.5,2.15)..(0.65,2.05)..controls (0.8,2.15) and (1.5,2.2)..(1.55,2.4);\draw (0.65,2.05)--(0.66,1.87);
\draw (2.4,1.1)..controls (2.2,0.75) and (2.3,0.65)..(3,0.62)..controls (3.1,0.4) and (3.1,0.4)..(3.1,0.32);\draw (3,0.62)..controls (3.6,0.65) and (4,0.65)..(4.1,0.9);
\draw[->] (0.4,0)--(0.85,0.55); \draw[->] (4.2,-0.2)--(4.65,0.4);
\end{tikzpicture}}
\mspage{283}
